\documentclass[conference]{IEEEtran}

\usepackage{amsmath,amssymb,amsfonts}
\usepackage{graphicx}
\usepackage{textcomp}
\usepackage{xcolor}
\usepackage{booktabs}
\usepackage{array}
\usepackage{listings}
\usepackage{url}
\usepackage{hyperref}
\usepackage{cite}

% Code Listing Style
\lstset{
	language=Python,
	basicstyle=\ttfamily\footnotesize,
	keywordstyle=\color{blue}\bfseries,
	commentstyle=\color{gray}\itshape,
	stringstyle=\color{red},
	numbers=left,
	numberstyle=\tiny\color{gray},
	numbersep=5pt,
	frame=single,
	breaklines=true,
	captionpos=b,
	tabsize=4
}

\def\BibTeX{{\rm B\kern-.05em{\sc i\kern-.025em b}\kern-.08em
		T\kern-.1667em\lower.7ex\hbox{E}\kern-.125emX}}

\begin{document}
	
	\title{[Your Project Title Here]\\
		{\footnotesize BIL 476 / 573 --- Data Mining Course Project, Summer 2026\\
			TOBB University of Economics and Technology}
	}
	
	\author{\IEEEauthorblockN{[Student Full Name]}
		\IEEEauthorblockA{\textit{TOBB University of Economics and Technology}\\
			Ankara, Turkey \\
			Student ID: [XXXXXXXXX] \\
			student@etu.edu.tr}
	}
	
	\maketitle
	
	\begin{abstract}
		Provide a concise summary of your project (150--250 words). State the problem you addressed, the data mining technique(s) used, the dataset, key results, and main conclusions. Do not use symbols, special characters, footnotes, or math in the abstract.
	\end{abstract}
	
	\begin{IEEEkeywords}
		data mining, [keyword 2], [keyword 3], [keyword 4], [keyword 5]
	\end{IEEEkeywords}
	
	\begin{center}
	\fbox{\begin{minipage}{0.92\linewidth}
	\textbf{Template note.} Replace all instructional text, examples, and placeholders before submission. The final report must contain your own analysis and must include the AI Tool Use and Originality Declaration before the References section.
	\end{minipage}}
	\end{center}
	
	\section{Introduction}
	\label{sec:introduction}
	
	Describe the problem you are addressing. Why is it important? What is the real-world motivation? Clearly state your research question or objective.
	
	\textbf{Guidelines:}
	\begin{itemize}
		\item Define the problem clearly and concisely.
		\item Explain why this problem is relevant and important.
		\item State your specific research question or objective.
		\item Briefly mention the approach you took.
		\item Outline the structure of the rest of the paper.
		\item Make sure the motivation, research question, and claimed contribution are your own, not copied from an AI-generated outline.
	\end{itemize}
	
	\textit{Example: ``In this project, we investigate whether frequent pattern mining can reveal meaningful purchasing behaviors in online retail data. We compare the effectiveness of Apriori and FP-Growth algorithms under different support thresholds and evaluate the discovered rules using multiple interestingness measures.''}
	
	\section{Related Work}
	\label{sec:related_work}
	
	Provide a literature review of relevant prior work.
	
	\textbf{Requirements:}
	\begin{itemize}
		\item \textbf{BIL 476 (Undergraduate):} Cite at least 3--5 references.
		\item \textbf{BIL 573 (Graduate):} Cite at least 5--8 references.
	\end{itemize}
	
	Discuss how your work relates to or differs from existing studies. For example: ``Han et al.~\cite{b1} introduced the FP-Growth algorithm as an efficient alternative to Apriori. Smith and Jones~\cite{b2} applied K-Means clustering to customer segmentation but did not consider the effect of dimensionality reduction, which we address in this work.''
	
	Do not cite sources that you have not read. If AI tools helped you discover papers or summarize candidate sources, you must verify the sources from the original publications and write the synthesis yourself.
	
	\section{Dataset Description}
	\label{sec:dataset}
	
	\subsection{Data Source}
	Describe where the data comes from (e.g., Kaggle, UCI Machine Learning Repository, data.gov, web scraping, synthetic generation). Include the URL or reference. Explain why this dataset is appropriate for your chosen problem.
	
	\subsection{Data Characteristics}
	Provide a summary of the dataset using the table format below:
	
	\begin{table}[htbp]
		\caption{Dataset Summary}
		\begin{center}
			\begin{tabular}{|l|l|}
				\hline
				\textbf{Property} & \textbf{Value} \\
				\hline
				Number of Records & e.g., 10,000 \\
				\hline
				Number of Attributes & e.g., 15 \\
				\hline
				Attribute Types & Nominal, Ordinal, Numeric \\
				\hline
				Target Variable & e.g., Class (Yes/No) \\
				\hline
				Missing Values & e.g., 5\% in column X \\
				\hline
				Class Distribution & e.g., 70\% / 30\% \\
				\hline
				Source URL & e.g., kaggle.com/datasets/... \\
				\hline
			\end{tabular}
			\label{tab:dataset}
		\end{center}
	\end{table}
	
	Include descriptive statistics (mean, median, standard deviation, min, max) for key numeric attributes. Show sample rows if helpful.
	
	\subsection{Ethical Considerations}
	Discuss any privacy, bias, or ethical concerns related to the data. How did you address them? Do not upload confidential, private, restricted, or non-public datasets to AI tools. If any AI tool was used while working with data, state what was shared and confirm that privacy, confidentiality, and licensing restrictions were respected.
	
	\section{Methodology}
	\label{sec:methodology}
	
	\subsection{Data Preprocessing}
	Describe the preprocessing steps applied and justify each choice:
	\begin{itemize}
		\item \textbf{Missing value handling:} Method used (e.g., mean/median imputation, deletion, model-based) and why.
		\item \textbf{Noise removal:} Smoothing, outlier treatment.
		\item \textbf{Normalization / Standardization:} e.g., min-max scaling, z-score normalization.
		\item \textbf{Feature selection / Dimensionality reduction:} e.g., PCA, correlation-based selection.
		\item \textbf{Data transformation:} e.g., discretization, encoding categorical variables.
	\end{itemize}
	
	Discuss the \textbf{effect of different preprocessing choices} on your results where applicable.
	
	\subsection{Algorithm Selection and Description}
	Describe the algorithm(s) you chose and \textbf{justify why} they are appropriate for your problem. Provide a brief theoretical background for each algorithm.
	
	You must apply \textbf{at least two different approaches or variations} and compare them. For example:
	\begin{itemize}
		\item Different algorithms for the same task (e.g., K-Means vs. DBSCAN)
		\item Same algorithm with different parameters or data transformations
		\item Standard method vs. a hybrid/novel approach (especially for BIL 573 students)
	\end{itemize}
	
	\textbf{BIL 573 (Graduate) students} are expected to go beyond standard methods and explore hybrid, modified, or novel approaches.
	
	\subsection{Experimental Setup}
	\begin{itemize}
		\item Train/test split ratio (e.g., 80/20) or cross-validation strategy (e.g., 10-fold)
		\item Hyperparameter settings and how they were chosen
		\item Hardware and software environment
		\item Programming language and libraries used (e.g., Python 3.x, scikit-learn, mlxtend, pandas, matplotlib)
	\end{itemize}
	

	\subsection{Reproducibility and Verification}
	State how another student or evaluator can reproduce your results. Provide the code repository, random seeds, software versions, and execution instructions where possible. If any AI-assisted code was used, you must review, test, and understand it. You remain responsible for every line of submitted code and for every numerical result reported in the paper.
	
	\section{Results}
	\label{sec:results}
	
	\subsection{Evaluation Metrics}
	Describe the metrics used and justify their appropriateness for your task. Examples by task type:
	\begin{itemize}
		\item \textbf{Classification:} Accuracy, Precision, Recall, F1-Score, ROC/AUC
		\item \textbf{Clustering:} Silhouette Score, SSE, Davies-Bouldin Index
		\item \textbf{Pattern Mining:} Support, Confidence, Lift, Kulczynski
		\item \textbf{Anomaly Detection:} Precision, Recall, F1, False Positive Rate
		\item \textbf{Regression/Prediction:} RMSE, MAE, $R^2$
	\end{itemize}
	
	\subsection{Results Summary}
	Present your results with appropriate tables and figures.
	
	\begin{table}[htbp]
		\caption{Classification Results Comparison (Example)}
		\begin{center}
			\begin{tabular}{|l|c|c|c|c|}
				\hline
				\textbf{Model} & \textbf{Accuracy} & \textbf{Precision} & \textbf{Recall} & \textbf{F1} \\
				\hline
				Decision Tree & 0.85 & 0.83 & 0.80 & 0.81 \\
				\hline
				Na\"ive Bayes & 0.82 & 0.79 & 0.84 & 0.81 \\
				\hline
				Random Forest & 0.89 & 0.87 & 0.86 & 0.86 \\
				\hline
			\end{tabular}
			\label{tab:results}
		\end{center}
	\end{table}
	
	\subsection{Visualizations}
	Include relevant figures such as:
	\begin{itemize}
		\item Confusion matrices, ROC curves (classification)
		\item Cluster scatter plots, dendrograms (clustering)
		\item Frequent itemset/rule visualizations (pattern mining)
		\item Time series plots (prediction)
		\item Boxplots, histograms, heatmaps (EDA)
	\end{itemize}
	
	Use the abbreviation ``Fig.~\ref{fig:example}'' when referencing figures, even at the beginning of a sentence.
	
	\begin{figure}[htbp]
		\centerline{\rule{0.35\textwidth}{3cm}}
		\caption{Replace with your actual figure. Use 8pt Times New Roman for axis labels. Write full quantity names (e.g., ``Accuracy (\%)''), not just symbols.}
		\label{fig:example}
	\end{figure}
	
	\subsection{Comparison of Approaches}
	Compare the different methods/variations you tried. Discuss:
	\begin{itemize}
		\item Which performed best and \textbf{why}?
		\item How did different data transformations or preprocessing choices affect results?
		\item Statistical significance of differences (if applicable).
	\end{itemize}
	
	\section{Discussion}
	\label{sec:discussion}
	
	Provide a critical analysis of your work:
	\begin{itemize}
		\item Which method performed best and \textbf{why}?
		\item What were the \textbf{limitations} of your approach?
		\item Were there any \textbf{surprising or unexpected findings}?
		\item How did data quality issues (missing values, noise, imbalance) affect results?
		\item Discuss any \textbf{ethical implications} (bias, fairness, privacy).
		\item \textbf{BIL 573 students:} Discuss the novelty of your approach and its potential for academic publication.
	\end{itemize}
	
	The discussion must reflect your own interpretation of the results. AI tools may be used for limited language polishing or for asking clarifying questions, but the argument, interpretation, limitations, and future-work directions must be written and justified by you.
	
	\section{Conclusion and Future Work}
	\label{sec:conclusion}
	
	Summarize your key findings concisely. Address:
	\begin{itemize}
		\item What did you achieve?
		\item What are the main takeaways?
		\item What would you do differently if you had more time?
		\item Suggest concrete directions for future work.
	\end{itemize}
	

	\section*{AI Tool Use and Originality Declaration}
	\label{sec:ai_declaration}
	
	This section is required in every submitted report. Keep the version that applies to your project and delete the other examples.
	
	\textbf{Example if no AI tools were used:} No generative AI or AI-assisted tools were used to generate the submitted report text, code, analysis, figures, interpretations, or references.
	
	\textbf{Example if limited AI support was used:} AI-assisted tools were used only for limited support, such as concept clarification, code debugging, code review, grammar improvement, and/or LaTeX formatting. The project idea, methodology choices, experiments, analysis, interpretation of results, and final report wording were reviewed, verified, and finalized by the student. The student accepts full responsibility for the accuracy, originality, reproducibility, and integrity of the submitted work.
	
	If AI tools were used, complete Table~\ref{tab:ai_use}. Use the 0--6 scale as follows: 0 = not used; 1 = minimal hints; 2 = low support with student-authored final content; 3 = regular support for specific sub-tasks with critical student review; 4 = high support for critical sub-tasks; 5 = very high AI contribution; 6 = AI-led work. Levels 5--6 may indicate that the work falls outside acceptable use for this project and may negatively affect the grade.
	
	\begin{table}[htbp]
		\caption{AI Use Disclosure Summary}
		\label{tab:ai_use}
		\centering
		\footnotesize
		\begin{tabular}{|p{0.35\linewidth}|c|p{0.38\linewidth}|}
			\hline
			\textbf{Project Stage} & \textbf{Level} & \textbf{Brief Description} \\
			\hline
			Conceptualization and idea generation & [0--6] & [Explain briefly] \\
			\hline
			Literature search and synthesis & [0--6] & [Explain briefly] \\
			\hline
			Methodology, coding, and analysis & [0--6] & [Explain briefly] \\
			\hline
			Writing, editing, and LaTeX formatting & [0--6] & [Explain briefly] \\
			\hline
			Interpretation, discussion, and conclusion & [0--6] & [Explain briefly] \\
			\hline
		\end{tabular}
	\end{table}
	
	Do not list an AI tool as an author or as a source of factual claims. AI-generated references, explanations, code, figures, or numerical results must be independently checked before inclusion. The instructor may ask you to explain your work orally, reproduce part of the analysis, or provide a brief record of prompts and outputs if the submitted work does not clearly demonstrate your own understanding.
	

	\begin{thebibliography}{00}
		
		\bibitem{b1} J.~Han, M.~Kamber, and J.~Pei, \textit{Data Mining: Concepts and Techniques}, 3rd~ed. Morgan Kaufmann, 2011.
		
		\bibitem{b2} A.~Smith and B.~Jones, ``Customer segmentation using K-Means clustering,'' in \textit{Proc. IEEE Int. Conf. Data Mining}, 2023, pp.~100--110.
		
	\end{thebibliography}
	
	\newpage	
	\appendix
	
	\section{Code Repository}
	\label{app:code}
	
	Full source code is available at: \url{https://github.com/yourusername/project-repo}
	
	\begin{lstlisting}[caption={Example: Data Loading and Preprocessing}, label={lst:preprocess}]
		import pandas as pd
		import numpy as np
		from sklearn.model_selection import train_test_split
		from sklearn.preprocessing import StandardScaler
		
		# Load dataset
		df = pd.read_csv('data/dataset.csv')
		print(f"Dataset shape: {df.shape}")
		print(f"Missing values:\n{df.isnull().sum()}")
		
		# Handle missing values
		df.fillna(df.median(numeric_only=True), inplace=True)
		
		# Split data
		X = df.drop('target', axis=1)
		y = df['target']
		X_train, X_test, y_train, y_test = train_test_split(
		X, y, test_size=0.2, random_state=42,
		stratify=y
		)
		
		# Normalize
		scaler = StandardScaler()
		X_train = scaler.fit_transform(X_train)
		X_test = scaler.transform(X_test)
	\end{lstlisting}
	
	\section{AI Prompt and Verification Log (If Applicable)}
	\label{app:ai_log}
	
	If AI tools were used, you may include a short table summarizing important prompts, the purpose of use, what output was accepted or rejected, and how the output was verified. Do not include private data, passwords, API keys, or confidential material.
	
	\section{Additional Figures and Tables}
	\label{app:additional}
	
	Include any supplementary material here (additional plots, extended result tables, hyperparameter tuning details, etc.).
	
\end{document}
